\RequirePackage[T1]{fontenc}
\documentclass[10pt,conference]{IEEEtran}
\usepackage{amsmath,amssymb,amsthm,graphicx,booktabs,array,url,flushend}
\usepackage[hidelinks]{hyperref}
\usepackage{tikz}
\usetikzlibrary{arrows.meta,positioning}
\emergencystretch=2em
\newtheorem{theorem}{Theorem}
\newtheorem{lemma}{Lemma}
\newtheorem{proposition}{Proposition}
\theoremstyle{definition}
\newtheorem{definition}{Definition}
\newtheorem{remarknum}{Remark}
\newcommand{\code}[1]{\texttt{#1}}
\newcommand{\G}[1]{G#1}
\newcommand{\tool}{ReleaseWeave}
\newcommand{\lang}{RIL}
\begin{document}
\title{Recovering Versioned Releases Across Build and Publish Boundaries}
\maketitle
\begin{abstract}
An npm package can build from its dependency's local workspace output, while
an isolated registry-based Docker builder needs its base image published
first. After partial failure, a consumer can also
lose its pending release when its provider's success removes the original
propagation reason. We present a protocol combining stage-aware scheduling
with consumer-based admission of release intent. Git history determines the
plan; tags record completed obligations. We formalize the single-history core
with deterministic planning and eight conditional recovery guarantees, and
describe its extension to polyrepo fleets with repository-owned records,
central coordination, or peer choreography. Fleet recovery requires evidence
of which source revisions each consumer incorporated.
All planned outcomes hold in a prospective 140-episode Docker matrix; a
source-window counterfactual loses 111 obligations in 84/120 fault episodes.
A 772-state finite projection finds no model violations and 260 losing states.
Companion evidence adds 300 reference-model agreements, 20 recovery
episodes, 60 mixed-chain executions, and three CI repetitions; historical
observations add broader diagnostics, comparisons, and scale. The guarantees
concern the model; implementation refinement and external exactly-once effects
remain outside the claims.
\end{abstract}
\begin{IEEEkeywords}
software release engineering, continuous integration and delivery,
dependency scheduling, recovery, idempotence
\end{IEEEkeywords}
\input{introduction}
\input{system}
\input{model}
\input{guarantees}
\input{evaluation}
\input{discussion}
\clearpage
\IEEEtriggeratref{34}
\input{references}
\end{document}
